\chapter*{Заключение}
\addcontentsline{toc}{chapter}{Заключение}

Выпускная квалификационная работа была посвящена исследованию и сравнительному анализу методов прогнозирования валютного курса USD/UZS на основе классических моделей временных рядов, алгоритмов машинного обучения и современных нейросетевых архитектур.

Актуальность выбранной темы обусловлена высокой практической значимостью прогнозирования валютных курсов для финансовых организаций, предприятий внешнеэкономической деятельности, государственных структур и частных инвесторов. В условиях нестабильности финансовой среды и мировой экономики возрастает острая потребность в инструментах, позволяющих получать максимально точные и обоснованные прогнозы изменения котировок.

На этапе теоретического исследования были систематизированы основы анализа финансовых временных рядов, а также проведено математическое обоснование применения традиционных статистических методов прогнозирования, алгоритмов машинного обучения и рекуррентных нейросетевых моделей. Особое внимание было уделено специфике валютного курса как сложного экономического показателя, который характеризуется нестационарностью, высоким уровнем информационного шума и поведением, близким к случайному блужданию.

В рамках практической реализации в качестве объекта исследования был выбран официальный курс доллара США к узбекскому суму (USD/UZS). Для проведения корректного анализа использовались исторические данные исключительно за период после либерализации валютного рынка Республики Узбекистан (начиная с сентября 2017 года), что позволило исключить нерыночные структурные сдвиги. На этапе подготовки данных были выполнены все необходимые методологические процедуры: очистка временного ряда, нормализация значений и формирование лаговых признаков для обучения моделей.

В ходе эксперимента были программно реализованы и протестированы четыре конфигурации прогнозирования: статистическая модель ARIMA в двух режимах (rolling forecast с переобучением и статический прогноз без переобучения), алгоритм ансамблевого машинного обучения Random Forest и рекуррентная нейросетевая модель долгой краткосрочной памяти (LSTM). Для объективной оценки качества прогнозирования применялся комплекс стандартных статистических метрик: средняя абсолютная ошибка (MAE), среднеквадратичная ошибка (MSE), корень из среднеквадратичной ошибки (RMSE) и средняя абсолютная ошибка в процентах (MAPE).

Основные результаты исследования показали, что все рассмотренные модели способны воспроизводить общую динамику изменения валютного курса. Наименьшие абсолютные ошибки были получены для модели ARIMA с rolling forecast (MAPE = 0,15\%), однако данный подход использует механизм постоянного переобучения, что даёт методологическое преимущество перед моделями с однократным обучением. При корректном сравнении в равных условиях (без переобучения) наилучшие результаты продемонстрировала нейросетевая модель LSTM (MAPE = 0,33\%), значительно превзошедшая как Random Forest (MAPE = 0,88\%), так и статическую ARIMA (MAPE = 1,89\%).

Проведённое исследование позволяет сделать два ключевых вывода. Во-первых, для оперативного краткосрочного прогнозирования с возможностью постоянного обновления модели наиболее эффективна ARIMA с rolling forecast. Во-вторых, при необходимости автономного прогнозирования без переобучения нейросетевая модель LSTM демонстрирует значительное превосходство над классическими методами.

Практическая значимость работы заключается в возможности использования полученных результатов и построенных моделей для краткосрочного прогнозирования валютных курсов. Предложенные алгоритмы могут применяться при разработке аналитических систем поддержки принятия решений в банковской и финансовой сфере.

Перспективы дальнейших исследований связаны с переходом к многомерному анализу, включением дополнительных макроэкономических факторов (инфляции, процентных ставок, платежного баланса), применением более сложных нейросетевых архитектур и использованием моделей ARCH/GARCH для более точного учёта волатильности финансовых рынков.

Цель выпускной квалификационной работы достигнута, а поставленные задачи выполнены в полном объёме.